\begin{helper}
Fix the split-outside-blocker replacement data, an already chosen finite
supported atom family \(\Omega\), and already chosen paired cells
\(C_\omega,C'_\omega\).  On this same atom family one may define the
replacement index sets
\[
  R_j=\{i:i<j\},\qquad T_j=\{i:j\le i<n\},
\]
and the core-included adjacent boundary events \(L_k(\omega)\) and
\(R_k(\omega)\).  For \(k<n\), \(L_k(\omega)\) is the source event in which
some \(x\in B_k\) is active, beats every active embedded neighbor in
\(X\cup\{r\}\), satisfies every retained noncurrent source-blocker test
\(i\in T_k\), \(i\ne k\), and is not killed by the current blocker \(z_k\).
The target event \(R_k(\omega)\) is the analogous event with the transported
core, every already replaced noncurrent private-blocker test
\(i\in R_{k+1}\), \(i\ne k\), and the current private blocker
\(\beta(x,z_k)\).  The construction also records a flat representative of
the sigma-type map \(\beta\) at the current incidence, agreeing with
\(\beta(x,z_k)\) for every \(x\in B_k\).
\end{helper}
\begin{proof}
No new atom family is chosen.  Define \(R_j\) to be the finite set of blocker
indices whose natural-number value is \(<j\), and define \(T_j\) to be the
finite set of indices whose value is at least \(j\).  These definitions are
made inside the finite type \(\operatorname{Fin} n\), so they apply to the
same \(n,z_i,B_i,\Omega,C_\omega,C'_\omega\) that were already supplied by
the supported common-cell package.

For \(k<n\), define \(L_k(\omega)\) by the displayed source formula: choose a
candidate \(x\in X\), require \(x\in B_k\), require \(x\) to be active, require
the embedded-neighbor comparison against every active \(y\in X\cup\{r\}\),
require all retained noncurrent tests for \(i\in T_k\), \(i\ne k\), and
finally require the negation of the current \(z_k\)-blocking event.  This is
a definition, so the asserted formula for \(L_k(\omega)\) follows
immediately from the definition.

For \(R_k(\omega)\), first define the flat representative \(\beta_k\) by
using the original sigma-type value \(\beta(x,z_k)\) when \(x\in X\) and
\(x z_k\in E(G)\), and by an arbitrary fallback value otherwise.  If
\(x\in B_k\), then the displayed definition of \(B_k\) supplies exactly
\(x\in X\) and \(x z_k\in E(G)\), while the outside-blocker list supplies
\(z_k\notin X\cup\{r\}\).  Hence \(\beta_k(x,z_k)=\beta(x,z_k)\) in the
sigma-type notation.  Define \(R_k(\omega)\) by the transported core formula,
the already replaced noncurrent private tests, and the current private test
for \(\beta_k(x,z_k)\).  Again this is a definition on the original
\(\Omega,C_\omega,C'_\omega\), so the target adjacent-event formula and the
agreement clause follow directly.
\end{proof}
